% Einheit 2 von 4 – Äquivalenzumformungen (bank/lineare-gleichungen/e2.jsonl, zusammenbau v0.1)
%% TODO Zweigzeile Teil 1: was hier gelernt wird (Satz in Schülersprache aus dem Einheitstitel) – Einheit 2
\einheitenkopf[e2]{Einheit 2 von 4 · Äquivalenzumformungen}
\zweigzeile{ab Kl. 7, je nach Buch bis Kl. 8 · P10}
\setcounter{aufgabe}{11}

%% TODO Titel: Ich-kann-Satz fehlt in der Bank (Platzhalter: Kettenname) – Nr. 12
%% TODO Anweisung: ein Satz über der ersten Teilaufgabe fehlt in der Bank – Nr. 12
\begin{aufgabe}{Vorzahl und Zeichen bei x lesen}
\begin{teile}
\teil $-5x = 30$ – was steht bei $x$? \leerfeld
\end{teile}
\end{aufgabe}

%% TODO Titel: Ich-kann-Satz fehlt in der Bank (Platzhalter: Kettenname) – Nr. 13
%% TODO Anweisung: ein Satz über der ersten Teilaufgabe fehlt in der Bank – Nr. 13
\begin{aufgabe}{Reihenfolge bestimmen}
\begin{teile}
\teil $3x + 7 = 22$ – welche Umformung zuerst? $\mid$ \leerfeld
\end{teile}
\end{aufgabe}

%% TODO Titel: Ich-kann-Satz fehlt in der Bank (Platzhalter: Kettenname) – Nr. 14
%% TODO Anweisung: ein Satz über der ersten Teilaufgabe fehlt in der Bank – Nr. 14
\begin{aufgabe}{Umformen}
%% TODO \rechenplatz{n}: Schrittzahl des Grundfalls fehlt in der Bank (nur bei fester Schreibform, 2.3 b)
\begin{teile}
\teil $x + 11 = 18$ – was kommt hinter den Strich? $\mid$ \leerfeld
\end{teile}
\begin{gleichungsraster}[2]
\gl{x + 8 = 15} &
\gl{6x = 54} \\
\gl{x + 12 = 5} &
\gl{\text{$3x + 5 = 26$ – mit Probe}} \\
\gl{-3x = 18} &
\gl{\frac{x}{4} = 6} \\
\end{gleichungsraster}
\begin{teile}
\teil Eine zweischrittige Gleichung mit der Lösung 6?
\end{teile}
\begin{gleichungsraster}[2]
\gl{\text{$3(x + 7) = 12$ – löse und mache die Probe. (P10 2020 OS)}} & \\
\end{gleichungsraster}
\end{aufgabe}

%% TODO Titel: Ich-kann-Satz fehlt in der Bank (Platzhalter: Kettenname) – Nr. 15
%% TODO Anweisung: ein Satz über der ersten Teilaufgabe fehlt in der Bank – Nr. 15
\begin{aufgabe}{Umformen – Fehler finden, Begründen}
\begin{teile}
\teil Jana rechnet: \rechnung{3x &= 15 &&\mid -3 \\ x &= 12} Finde den Fehler.
\teil Leo rechnet: \rechnung{x + 8 &= 14 \\ x &= 14 + 8 \\ x &= 22} Finde den Fehler.
\teil Sara rechnet: \rechnung{3x + 9 &= 21 &&\mid :3 \\ x + 9 &= 7} Finde den Fehler.
\teil Warum darf man bei $x + 4 = 10$ auf beiden Seiten 4 abziehen?
\end{teile}
\end{aufgabe}
